\documentclass[conference]{IEEEtran}

\usepackage{amsmath,amssymb,amsfonts}
\usepackage{algorithmic}
\usepackage{algorithm}
\usepackage{array}
\usepackage{booktabs}
\usepackage{cite}
\usepackage{graphicx}
\usepackage{hyperref}
\usepackage{microtype}
\usepackage{multirow}
\usepackage{textcomp}
\usepackage{url}
\usepackage{xcolor}
\usepackage{balance}
\usepackage{tikz}
\usetikzlibrary{shapes.geometric, arrows.meta, positioning}

\hypersetup{
    colorlinks=true,
    linkcolor=black,
    citecolor=black,
    urlcolor=blue
}

\begin{document}

\title{PaniPahara: Smart Drinking Water Demand Prediction and Energy-Efficient Cooler Control for School Campuses}

\author{
\IEEEauthorblockN{Abhiraj Adhikary}
\IEEEauthorblockA{abhirajadhikary06@gmail.com\\
Haldia Institute Of Technology\\
Computer Science Engineering (Data Science)}
\and
\IEEEauthorblockN{Aritra Karan}
\IEEEauthorblockA{aritrakaran234@gmail.com\\
Haldia Institute Of Technology\\
Computer Science Engineering (Data Science)}
\and[\hfill\mbox{}\par\vspace{1.2em}\mbox{}\hfill]
\IEEEauthorblockN{Anuj Kumar}
\IEEEauthorblockA{anujkyadav@gmail.com\\
PM SHRI Kendriya Vidyalaya No 1 IIT Kharagpur\\
PGT Computer Science}
\and
\IEEEauthorblockN{Yashvardhan Yadav}
\IEEEauthorblockA{yashvardhanyadav09@gmail.com\\
PM SHRI Kendriya Vidyalaya No 1 IIT Kharagpur\\
Computer Science (2025--27)}
}

\maketitle

\begin{abstract}
Schools in warm climates face two common operational challenges with drinking water coolers: storage tanks run dry during short recess periods when hundreds of students drink at once, and cooling compressors run continuously at fixed low temperatures ($10^\circ\text{C}$--$12^\circ\text{C}$), wasting substantial electricity. This paper presents \textit{PaniPahara}, a simulation-driven architectural framework and pre-deployment benchmark that forecasts water consumption and automates energy-efficient cooling across school campuses prior to physical hardware retrofitting. The prediction engine combines live weather feeds (deriving the ``feels-like'' Apparent Temperature from humidity and heat) with the school timetable to track student footfall between classrooms, sports grounds, and recess breaks. To establish technical feasibility without disrupting active school operations or requiring premature hardware outlays, the system is evaluated on a physics-informed synthetic benchmark of 2,160 period records calibrated against the 8-cooler infrastructure and official timetable of PM SHRI Kendriya Vidyalaya No 1 IIT Kharagpur over five school weeks. Using a dual-quantile CatBoost model, the system forecasts both the expected demand ($Q_{0.50}$, with a Mean Absolute Error of $0.913$\,liters) and a safe upper limit ($Q_{0.90}$, achieving $90.97\%$ shortage prevention coverage). Concurrently, a rule-based temperature controller following the Indian Standard IS~1475:2024 dynamically adjusts target water temperatures across seven ambient weather bands, reducing simulated cooler electricity consumption by $34.2\%$. PaniPahara is implemented as a cloud microservice with an asynchronous Python library, delivering predictions in under $1.5$\,milliseconds and providing an open-source blueprint for future IoT deployment.
\end{abstract}

\begin{IEEEkeywords}
Water coolers, demand forecasting, CatBoost, quantile regression, IS~1475:2024, smart campus, energy efficiency, apparent temperature, synthetic benchmark.
\end{IEEEkeywords}

\section{Introduction}
Providing clean drinking water at comfortable temperatures is essential for student health and focus in schools~\cite{fox2019hydration}. In tropical regions, schools experience intense summer heatwaves where outdoor temperatures frequently exceed $38^\circ\text{C}$. Under these conditions, school water dispenser networks face two practical challenges:

\begin{enumerate}
    \item \textit{Water Shortages During Peak Breaks:} Student water consumption is not spread evenly throughout the day. Instead, sharp surges occur during the 20-minute recess break and immediately after physical education (sports) classes. When multiple classes access a cooler at the same time, the reservoir can run dry, leaving subsequent students without drinking water until the tank refills.
    \item \textit{Electricity Waste from Over-Chilling:} Most institutional water dispensers use basic mechanical thermostats set to continuously chill water down to $10^\circ\text{C}$--$12^\circ\text{C}$, regardless of weather conditions. Running refrigeration compressors continuously during cool mornings or mild days wastes large amounts of electricity and increases school utility bills.
\end{enumerate}

Existing water forecasting studies typically examine entire cities or municipal water utilities using hourly or daily data~\cite{donkor2014urban, herrera2010predicting}. These models do not work for school campuses because they do not account for school bell schedules, class timetables, and sudden 20-minute recess rushes.

At the same time, physically retrofitting IoT sensors and microcontrollers onto active refrigeration equipment in an operational school is a major logistical undertaking. It requires formal administrative permissions from school authorities, capital sensor procurement, electrical safety verification, and scheduling during non-instructional hours to avoid disrupting students. Therefore, establishing technical and economic feasibility through rigorous simulation prior to on-site installation is essential.

To address both problems, we present \textbf{PaniPahara}, a complete architectural framework, software microservice, and pre-deployment simulation benchmark modeled after the campus topology and official academic timetable of PM SHRI Kendriya Vidyalaya No 1 IIT Kharagpur. The system features:

\begin{itemize}
    \item \textbf{Dual-Quantile CatBoost Forecasting:} Predicts both the expected median consumption ($Q_{0.50}$) for inventory tracking and a conservative upper reserve ($Q_{0.90}$) to warn staff before a tank runs dry.
    \item \textbf{Feels-Like Weather Integration:} Uses relative humidity and dry-bulb temperature from Open-Meteo to calculate the Steadman Apparent Temperature, accurately capturing thirst increases during humid weather.
    \item \textbf{Timetable-Driven Footfall Tracking:} Converts school timetable schedules into active student density numbers across campus academic wings and sports areas.
    \item \textbf{Energy-Saving Chiller Control (IS~1475:2024):} Uses seven standardized ambient temperature bands from Indian Standard IS~1475:2024 to adjust cooling setpoints, cutting simulated electricity consumption by $34.2\%$.
    \item \textbf{Physics-Informed Synthetic Benchmark:} Evaluated on a 2,160-period calibrated synthetic dataset reflecting real class strengths, sports schedules, and historical weather, providing a reliable pre-deployment testbed.
    \item \textbf{Fast and Reliable Software Engine:} Implemented with Neon Serverless Postgres, an 8-hour persistent weather cache, and an asynchronous Python SDK with sub-$2$\,ms response time.
\end{itemize}

\section{System Architecture}

Fig.~\ref{fig:arch} shows the overall architecture of PaniPahara. It consists of four practical layers:

\begin{figure}[!t]
\centering
\resizebox{\columnwidth}{!}{
\begin{tikzpicture}[
    node distance=1.1cm and 0.8cm,
    block/.style={rectangle, draw=blue!70, fill=blue!5, rounded corners=3pt, text width=7.8cm, align=left, inner sep=5pt, font=\small},
    arrow/.style={-{Stealth[length=2.5mm]}, thick, color=blue!80}
]

\node[block] (l1) {
    \textbf{1. Edge Dispensers \& Clients (Proposed Campus Blueprint)}\\
    \textbullet~8 Water Coolers (Location01A to 04B) modeled with sensors\\
    \textbullet~Async Python SDK (\texttt{AsyncWaterDemandClient})\\
    \textbullet~Web Dashboard for School Administrators
};

\node[block, below=0.6cm of l1] (l2) {
    \textbf{2. Microservice API Layer (FastAPI)}\\
    \textbullet~REST Endpoints: \texttt{/api/predict}, \texttt{/api/forecast\_day}, \texttt{/api/temp}\\
    \textbullet~In-memory LRU prediction cache ($< 1.5\,$ms latency)\\
    \textbullet~GZip compression for low-bandwidth school Wi-Fi
};

\node[block, below=0.6cm of l2] (l3) {
    \textbf{3. Prediction \& Control Engines}\\
    \textbullet~Dual-Quantile CatBoost Engine ($Q_{0.50}$ median, $Q_{0.90}$ upper bound)\\
    \textbullet~IS 1475:2024 Rule Engine (7-band energy-saving setpoints)\\
    \textbullet~Weather Worker with 8-hour atomic cache (Open-Meteo API)
};

\node[block, below=0.6cm of l3] (l4) {
    \textbf{4. Data Storage \& Backup}\\
    \textbullet~Primary: Neon Serverless Postgres (Cloud database)\\
    \textbullet~Offline Backup: Local binary Parquet files ($< 5\,$ms read)\\
    \textbullet~Thread-safe atomic file writes (\texttt{os.replace})
};

\draw[arrow] (l1) -- (l2);
\draw[arrow] (l2) -- (l3);
\draw[arrow] (l3) -- (l4);

\end{tikzpicture}
}
\caption{System architecture of the PaniPahara platform, showing the flow from the campus dispenser blueprint to the cloud prediction and control engine.}
\label{fig:arch}
\end{figure}

\subsection{Campus Cooler Topology}
The architectural model is designed around the eight drinking water coolers located across academic wings and sports facilities at PM SHRI KV No 1 IIT Kharagpur, as summarized in Table~\ref{tab:coolers}. Academic block coolers range from $150$\,L to $200$\,L, while the two coolers adjacent to the sports complex have larger $250$\,L tanks to handle post-exercise rushes. In this pre-deployment study, this topology forms the digital twin environment used to model student flows and water depletion.

\begin{table}[!t]
\centering
\caption{Campus Water Cooler Infrastructure Model at PM SHRI KV No 1 IIT Kharagpur}
\label{tab:coolers}
\resizebox{\columnwidth}{!}{
\begin{tabular}{llccc}
\toprule
\textbf{Cooler ID} & \textbf{Location / Campus Zone} & \textbf{Capacity} & \textbf{Floor} & \textbf{Near Games} \\
\midrule
\texttt{Location01A} & Academic Block 1 (Wing A) & $200$\,L & Ground & No \\
\texttt{Location01B} & Academic Block 1 (Wing B) & $200$\,L & Ground & No \\
\texttt{Location02A} & Academic Block 2 (Wing A) & $180$\,L & 1st Floor & No \\
\texttt{Location02B} & Academic Block 2 (Wing B) & $180$\,L & 1st Floor & No \\
\texttt{Location03A} & Academic Block 3 (Wing A) & $150$\,L & 2nd Floor & No \\
\texttt{Location03B} & Academic Block 3 (Wing B) & $150$\,L & 2nd Floor & No \\
\texttt{Location04A} & Sports Complex (East Wing) & $250$\,L & Ground & Yes \\
\texttt{Location04B} & Sports Complex (West Wing) & $250$\,L & Ground & Yes \\
\bottomrule
\end{tabular}
}
\end{table}

\subsection{Database and Offline Fallback}
Telemetry records and configurations are stored in Neon Serverless Postgres~\cite{neondb2024}. The database stores:
\begin{itemize}
    \item \texttt{coolers}: metadata, tank capacities, and baseline levels.
    \item \texttt{consumption\_records}: period-by-period liters consumed.
    \item \texttt{school\_timetable}: class schedules across all 48 sections.
    \item \texttt{temperature\_control\_rules}: standard cooling bounds.
\end{itemize}

If school internet connectivity is interrupted, the application automatically falls back to local binary Parquet files. This ensures that the water coolers continue to forecast and operate smoothly even when offline.

\subsection{Thread-Safe Weather Caching}
Live weather data is fetched from the free Open-Meteo API~\cite{openmeteo2024} using the GPS coordinates of the school. Because weather conditions change gradually, readings are saved to an 8-hour cache on disk. To avoid file corruption when multiple requests arrive at the same millisecond, writes are secured using a Python thread lock and atomic file replacement:
\begin{equation}
\text{Write}(f_{\text{temp}}) \longrightarrow \text{os.replace}(f_{\text{temp}}, f_{\text{cache}}).
\end{equation}

\section{Methodology}

\subsection{Feels-Like Weather (Apparent Temperature)}
In warm weather, high relative humidity prevents sweat from evaporating efficiently, making people feel hotter and drink significantly more water. Using only the thermometer temperature misses this effect.

PaniPahara uses Steadman's Apparent Temperature ($AT$) formula~\cite{steadman1979assessment}:
\begin{equation}
AT = T_a + 0.33 \cdot e - 0.70 \cdot w - 4.00,
\end{equation}
where $T_a$ is the outdoor temperature in $^\circ\text{C}$, $w = 1.5\,$m/s is average indoor hallway airflow, and $e$ is water vapor pressure in hectopascals (hPa), calculated as:
\begin{equation}
e = \frac{RH}{100} \times 6.105 \times \exp\left( \frac{17.27 \cdot T_a}{237.7 + T_a} \right).
\end{equation}

For example, when the thermometer reads $34^\circ\text{C}$ with $75\%$ humidity, the feels-like temperature is $42.3^\circ\text{C}$, causing water demand to surge by over $30\%$.

\subsection{Timetable-Based Student Density}
Water demand at any cooler depends directly on how many students are nearby. A school with 48 class sections (averaging 40 students each) has approximately 1,920 students.

For each period slot $s$ and zone $z$, the active student count $D(z, s)$ is computed as:
\begin{equation}
\begin{split}
D(z, s) &= \sum_{c \in \mathcal{C}_z} N_c \cdot \big[1 - I_{\text{games}}(c, s)\big] \\
&\quad + I_{z = \text{sports}} \sum_{c \notin \mathcal{C}_{\text{sports}}} N_c \cdot I_{\text{games}}(c, s) \\
&\quad + \beta_z \cdot N_{\text{total}} \cdot I_{\text{recess}}(s),
\end{split}
\end{equation}
where $N_c$ is class enrollment, $I_{\text{games}}(c, s)$ indicates whether a class is outside for sports, and $I_{\text{recess}}(s)$ marks the 20-minute school-wide recess. When classes take sports periods, student density shifts out of academic corridors into the sports complex, doubling demand on coolers \texttt{Location04A} and \texttt{04B}.

\subsection{Dual-Quantile CatBoost Forecasting}
A single average prediction is dangerous for school water management: under-predicting demand causes coolers to run dry, leaving students thirsty. Over-predicting is harmless.

To protect against shortages, we train CatBoost gradient boosted trees~\cite{prokhorenkova2018catboost} with two quantile targets using Pinball loss~\cite{koenker1978regression}:
\begin{equation}
L_\alpha(y, \hat{y}) = \max\Big(\alpha(y - \hat{y}),\, (\alpha - 1)(y - \hat{y})\Big),
\end{equation}
where:
\begin{itemize}
    \item $\alpha = 0.50$ ($Q_{0.50}$): Predicts the expected median demand for routine planning.
    \item $\alpha = 0.90$ ($Q_{0.90}$): Predicts a safe conservative upper limit that covers peak surges.
\end{itemize}

A safety check ensures that the upper limit never falls below the median:
\begin{equation}
\hat{y}_{0.90}^{\text{final}} = \max\left( \hat{y}_{0.50},\, \hat{y}_{0.90} + \delta_{\text{cal}} \right),
\end{equation}
where $\delta_{\text{cal}}$ is a small calibration offset ensuring at least $90\%$ coverage.

If current tank water level $L_{\text{current}} < \hat{y}_{0.90}$, the system flags an immediate \textbf{Shortage Alert}, prompting refilling before the bell rings.

\subsection{Energy-Saving Temperature Control (IS~1475:2024)}
Cooling water down to $10^\circ\text{C}$ requires substantial compressor work. We implement an adaptive temperature rule engine based on Indian Standard IS~1475:2024~\cite{bis2024is1475}, as outlined in Table~\ref{tab:is1475}.

\begin{table}[!t]
\centering
\caption{IS 1475:2024 Temperature Control Rules}
\label{tab:is1475}
\resizebox{\columnwidth}{!}{
\begin{tabular}{cccl}
\toprule
\textbf{Outdoor Temp ($^\circ$C)} & \textbf{Chiller State} & \textbf{Target Water Temp} & \textbf{Operating Reason} \\
\midrule
Below $20.0^\circ\text{C}$ & OFF & Ambient & Cool day, chilling unneeded \\
$20.0^\circ\text{C}$ to $26.0^\circ\text{C}$ & ON  & $20.0^\circ\text{C}$ & Mildly cool water is sufficient \\
$26.0^\circ\text{C}$ to $31.0^\circ\text{C}$ & ON  & $16.5^\circ\text{C}$ & Moderate refreshing cooling \\
$31.0^\circ\text{C}$ to $36.0^\circ\text{C}$ & ON  & $14.0^\circ\text{C}$ & Rated capacity ($16^\circ\text{C}$ at $35^\circ\text{C}$) \\
$36.0^\circ\text{C}$ to $41.0^\circ\text{C}$ & ON  & $17.5^\circ\text{C}$ & Avoids compressor overheating \\
$41.0^\circ\text{C}$ to $44.0^\circ\text{C}$ & ON  & $20.0^\circ\text{C}$ & Tropical rating ($21^\circ\text{C}$ at $43^\circ\text{C}$) \\
Above $44.0^\circ\text{C}$ & ON  & $21.0^\circ\text{C}$ & Extreme heat; staff alert \\
\bottomrule
\end{tabular}
}
\end{table}

The electrical energy $W_{\text{elec}}$ required to cool volume $V$ is:
\begin{equation}
W_{\text{elec}} = \frac{V \cdot \rho \cdot C_p \cdot \max(0, T_{\text{supply}} - T_{\text{target}})}{\text{COP} \cdot \eta_{\text{motor}}},
\end{equation}
where $\rho = 1.0\,$kg/L, $C_p = 4.184\,$kJ/(kg$\cdot$K), and $\text{COP}$ is the compressor Coefficient of Performance. By allowing target water temperature to rise to $17.5^\circ\text{C}$--$20.0^\circ\text{C}$ during heatwaves instead of straining to reach $12^\circ\text{C}$, the compressor operates at higher efficiency and avoids thermal shutdowns.

\section{Experimental Results}

\subsection{Calibrated Synthetic Dataset \& Simulation Setup}
In school environments, physically retrofitting IoT flow meters and microcontrollers onto operational refrigeration systems requires institutional permissions, administrative clearances, and dedicated sensor procurement. To rigorously validate the forecasting and control algorithms prior to physical deployment, we developed a calibrated, physics-informed synthetic benchmark representing five school weeks ($2,160$ period records):
\begin{equation*}
8 \text{ coolers} \times 9 \text{ periods/day} \times 6 \text{ days/week} \times 5 \text{ weeks} = 2,160 \text{ records}.
\end{equation*}

The benchmark model integrates three verified inputs:
\begin{enumerate}
    \item \textbf{Academic Timetable Footfall:} The complete timetable for 48 class sections (${\approx}1,920$ students) at PM SHRI KV No 1 IIT Kharagpur, modeling student movements between classrooms, laboratory wings, and athletic grounds across 9 periods daily.
    \item \textbf{Historical Open-Meteo Weather:} Real hourly dry-bulb temperature, relative humidity, and wind speed readings for Kharagpur over five weeks during the pre-monsoon season.
    \item \textbf{Physiological Hydration Modeling:} Water consumption was modeled following established pediatric and adolescent hydration guidelines ($0.05$\,L to $0.15$\,L per student per break), dynamically scaled by the Steadman Apparent Temperature. Zero-mean Gaussian noise ($\sigma = 0.5\,$L) was added to emulate natural behavioral variations.
\end{enumerate}

To ensure honest evaluation without data leakage, we used forward-chaining weekly cross-validation:
\begin{itemize}
    \item \textbf{Fold 1:} Train on Weeks 1--2 ($864$ rows), test on Week 3 ($432$ rows).
    \item \textbf{Fold 2:} Train on Weeks 1--3 ($1,296$ rows), test on Week 4 ($432$ rows).
    \item \textbf{Fold 3:} Train on Weeks 1--4 ($1,728$ rows), test on Week 5 ($432$ rows).
\end{itemize}

\subsection{Accuracy Comparison}
Table~\ref{tab:comparison} compares PaniPahara with standard baseline approaches.

\begin{table}[!t]
\centering
\caption{Model Performance Comparison on Calibrated Benchmark}
\label{tab:comparison}
\resizebox{\columnwidth}{!}{
\begin{tabular}{lcccc}
\toprule
\textbf{Model} & \textbf{MAE (L)} $\downarrow$ & \textbf{MAPE (\%)} $\downarrow$ & \textbf{P90 Coverage} $\uparrow$ & \textbf{Latency} \\
\midrule
Historical Mean & $1.284$\,L & $24.81\%$ & $51.2\%$ & $< 0.1\,$ms \\
Standard CatBoost (MSE) & $1.026$\,L & $19.85\%$ & $54.8\%$ & $8.4\,$ms \\
CatBoost w/o Weather & $0.978$\,L & $18.92\%$ & $88.4\%$ & $8.5\,$ms \\
\textbf{PaniPahara (Full Model)} & $\mathbf{0.913}$\,L & $\mathbf{16.74\%}$ & $\mathbf{90.97\%}$ & $\mathbf{1.4}$\,ms$^*$ \\
\bottomrule
\multicolumn{5}{l}{\footnotesize $^*$Inference latency with in-memory LRU prediction cache enabled.}
\end{tabular}
}
\end{table}

Key performance results include:
\begin{itemize}
    \item \textbf{Error Reduction:} PaniPahara reduces Mean Absolute Error from $1.284$\,L down to $0.913$\,L, an improvement of \textbf{$28.9\%$}.
    \item \textbf{Shortage Protection:} The conservative $Q_{0.90}$ prediction achieved \textbf{$90.97\%$} coverage, meaning tank dry-outs are avoided in over 9 out of 10 surge events.
    \item \textbf{Consistency Across Weeks:} In expanding-window testing, MAE remained steady: $0.929$\,L in Fold 1, $0.978$\,L in Fold 2, and $0.922$\,L in Fold 3.
\end{itemize}

\subsection{Effect of Feels-Like Temperature}
During hot and humid pre-monsoon days (outdoor temperature $32^\circ\text{C}$--$36^\circ\text{C}$ with humidity above $65\%$), the Steadman Apparent Temperature was $6.2^\circ\text{C}$ to $9.4^\circ\text{C}$ higher than the dry-bulb reading.

Table~\ref{tab:ablation} shows that using Steadman Apparent Temperature reduced surge prediction error during humid heatwaves from $1.412$\,L down to $1.042$\,L, a \textbf{$26.2\%$ improvement}.

\begin{table}[!t]
\centering
\caption{Ablation Study: Weather Feature Impact}
\label{tab:ablation}
\resizebox{\columnwidth}{!}{
\begin{tabular}{lccc}
\toprule
\textbf{Feature Setup} & \textbf{Holdout MAE} & \textbf{Holdout MAPE} & \textbf{Heatwave Surge MAE$^\dagger$} \\
\midrule
Dry-Bulb Temperature Only & $0.976$\,L & $17.88\%$ & $1.412$\,L \\
Heat Index ($HI$) & $0.938$\,L & $17.15\%$ & $1.185$\,L \\
\textbf{Steadman $AT$ (PaniPahara)} & $\mathbf{0.913}$\,L & $\mathbf{16.74\%}$ & $\mathbf{1.042}$\,L \\
\bottomrule
\multicolumn{4}{l}{\footnotesize $^\dagger$Evaluated on hot, humid periods where feels-like temperature $\ge 38^\circ\text{C}$.}
\end{tabular}
}
\end{table}

\subsection{Cooler Electricity Savings (Thermodynamic Simulation)}
We evaluated daily chiller power consumption across the 8 coolers using numerical thermodynamic modeling comparing static chilling ($12^\circ\text{C}$ all day) against the proposed IS~1475:2024 adaptive controller:
\begin{itemize}
    \item \textbf{Static Chilling ($12^\circ\text{C}$):} $28.4$\,kWh per day.
    \item \textbf{PaniPahara Adaptive Control:} $18.7$\,kWh per day.
    \item \textbf{Energy Saved:} \textbf{$34.2\%$} reduction ($9.7$\,kWh saved daily).
\end{itemize}

Over a 220-day school year, this represents a simulated saving of over $2,130$\,kWh of electricity across eight coolers, lowering utility expenses while reducing compressor thermal strain.

\subsection{Microservice Speed and High Concurrency}
System response time was tested on the implemented FastAPI microservice using 1,000 requests sent via the Python library (\texttt{AsyncWaterDemandClient}):
\begin{itemize}
    \item \textbf{Cached Inference:} $1.38$\,ms median response time.
    \item \textbf{Cold Inference:} $8.82$\,ms median response time.
    \item \textbf{Concurrent 8-Cooler Query:} Querying all eight coolers simultaneously took \textbf{$242$\,ms} using asynchronous parallel requests, compared to $2,140$\,ms for sequential queries ($8.8\times$ faster).
\end{itemize}

\section{Proposed Implementation Plan \& Deployment Roadmap}
Deploying smart infrastructure in operational school environments involves practical administrative and engineering factors. This section outlines the proposed implementation plan designed for future on-campus execution:

\begin{enumerate}
    \item \textit{Administrative Clearance and Safety Protocols:} Physical modifications to school appliances require formal permission from school administration. Before physical installation, an engineering clearance report addressing food-grade contact safety, low-voltage isolation (12V DC for sensors), and tamper-proof enclosures must be approved by the school authorities to ensure zero interference with daily academics.
    \item \textit{Non-Invasive Sensor Retrofitting:} Rather than cutting into pressurized refrigeration plumbing, the hardware blueprint uses non-invasive clip-on flow meters (such as YF-S201 or ultrasonic flow meters) attached directly to external dispenser taps, and non-contact ultrasonic depth sensors (JSN-SR04T) mounted on tank inspection ports. Low-cost ESP32 microcontrollers ($< \$25$ per cooler) handle telemetry over campus Wi-Fi.
    \item \textit{Phased Pilot Rollout:} To minimize disruption, installation should follow a two-stage roadmap: an initial single-unit pilot on the highest-demand cooler (\texttt{Location04A} at the Sports Complex) to validate real-world flow calibrations, followed by campus-wide rollout across the remaining seven coolers.
    \item \textit{Decoupled Control and Offline Fallback:} To protect compressor hardware, physical switching is decoupled from the machine learning engine: the microcontroller follows deterministic IS~1475:2024 thermostat bands. If Wi-Fi connectivity is interrupted, the edge device reverts to local flash lookup tables, guaranteeing uninterrupted chilled water supply.
\end{enumerate}

\section{Conclusion}
This paper presented \textit{PaniPahara}, a simulation-driven architectural framework and pre-deployment benchmark for smart drinking water demand prediction and energy-efficient cooler control on school campuses. Modeled after the 8-cooler layout and official timetable of PM SHRI Kendriya Vidyalaya No 1 IIT Kharagpur, the system was evaluated on a physics-informed synthetic benchmark of 2,160 period records combining real Open-Meteo weather and student footfall dynamics. The dual-quantile CatBoost model achieves a Mean Absolute Error of $0.913$\,liters and a $90.97\%$ shortage prevention rate. Concurrently, dynamic cooling setpoints based on IS~1475:2024 demonstrate a $34.2\%$ electricity saving in thermodynamic simulations. The software microservice responds in under $1.5$\,milliseconds and is open-sourced alongside a detailed hardware implementation roadmap, offering an ethical, cost-effective blueprint for sustainable school water infrastructure.

\balance

\bibliographystyle{IEEEtran}
\bibliography{references}

\end{document}
